\section{Reevaluation}

Each evaluation DAG represents an orbit's history which is valid 
%for 
with regards to the initial evaluation. %
%At reevaluation
When reevaluating, %the 
editing %of 
the parametric specification makes the topological parameters sub\-ject to changes. %
Thus, it is necessary to build reevaluation DAGs from the evaluation DAGs in order to update topological pa\-ra\-me\-ters. %
Once built, a reevaluation DAG can designate one, %or more 
several, 
or no orbit depending on the editing of the parametric specification. 
% \david{ou aucun si l'orbite est supprimée} %

In this section, we keep using the previous example (Fig.~\ref{fig:evaluation} and \ref{fig:reevaluation}) and its edited parametric specification (Fig.~\ref{fig:init-spec} and \ref{fig:reval-spec}), containing an added vertex insertion on the square's edge right after its creation and the deletion of the vertex insertion on the cube's edge. %

\subsection{Reevaluation DAG}

\begin{figure}[ht]
    \centering
    \includegraphics[height=.35\textheight]{pn3-reval}
    \caption{$PN_3$'s reevaluation DAG}
    \label{fig:pn3-reval-dag}
\end{figure}

Contrary to %its 
the related 
evaluation DAG, a reevaluation DAG is built top-down throughout the reevaluation process. %
%\david{
While an evaluation DAG represents the orbit's history of a topological parameter, the reevaluation DAG represents the history of this very same orbit after editing the parametric specification.
%}
%It describes the history of an orbit in the context of the new modeling process after edition. % While an evaluation DAG represents a topological parameter through the set of darts matched by rule nodes of the DAG, the reevaluation DAG represents the new designated topological parameters through the chosen darts. %

For example, let us consider $PN_{3}$ being the topological parameter of \texttt{4-extrusion}, which extrudes the cube's top face 
to produce %into 
a second cube. %
The $PN_3$'s reevaluation DAG (Fig.~\ref{fig:pn3-reval-dag}) is built step-by-step during the reevaluation process from its evaluation DAG (Fig.~\ref{fig:eval-dag-pn3}). %
\begin{description}
  \item[Level 1] The application \texttt{1-square} has no topological parameter and, thus, is identically reevaluated. Once again, its node $n5$ creates a single dart $5$ and the volume face $\orbit{0,1}(5)$ is identically reevaluated. This is why the orbit level 1 of the reevaluation DAG  references dart $5$. 
  \item[Level 2] The second application is an added one. \texttt{ADD1-insert} does match an edge of the tracked volume face $\orbit{0,1}(5)$ and modifies it, as deduced from the rule. Therefore, the event level 2 contains \texttt{MODIFICATION} and the orbit level 2 contains the same dart $5$ and orbit $\orbit{0,1}$.
  \item[Level 3] The third application is the extrusion of the square face into a cube. During its application, dart $5$ is matched by the hook $n_1$ (Fig.~\ref{fig:rule-face-extrusion}). From the evaluation DAG, the tracked dart is the copy of dart $5$ created by node $n_6$.
%  The rule's application 
Applying the rule
allows finding out this dart,
  numbered %is 
  $59$.
  \item[Level 4] Finally, the last application of $PN_3$ is deleted. Consequently, the modification that occurred during the initial e\-va\-lua\-tion does not occur at reevaluation. Therefore, the event level 4 contains \texttt{NOMODIF} and the orbit level 4 contains the same dart~$59$.
\end{description}
The reevaluation DAG %designates 
identifies 
the concrete name %from 
using 
its persistent name. 
%In the example above, 
$PN_3$'s concrete name is $59$ (cf. Fig.~\ref{fig:2-reeval} and \ref{fig:4-reeval}).


%Since darts can have different concrete names depending from the parametric specification editing, they must be found using parameter matching mechanism. % 

\begin{figure}[ht]
    \centering
    \includegraphics[height=.5\textheight]{pn6-reval}
    \caption{$PN_6$ reevaluation DAG}
    \label{fig:pn6-reval-dag}
\end{figure}

%This time, let's consider 
We now study 
a more complex example %being 
with the reevaluation DAG of $PN_{6}$ (Fig.~\ref{fig:pn6-reval-dag}),
which represents the pyramid's top vertex (Fig.~\ref{fig:6-eval}). %
\begin{description}
  \item[Level 1 to 4] 
  These three levels are similar to $PN_3$'s reevaluation DAG (Fig.~\ref{fig:pn3-reval-dag}), with the tracking of the volume face vertex $\orbit{1}$-orbit in addition to the $\orbit{0,1}$-orbit one. %
  The $\orbit{1}(5)$-orbit is crea\-ted by the \texttt{1-square} application. %
  \texttt{ADD1-insert} matches and preserves one dart of this orbit with the node $n_{0}$ of the insertion rule (Fig.~\ref{fig:rule-insert-vertex}).
  Thus, the event level 2 contains \texttt{NOMODIF}.
  Then, this orbit is copied to create a new volume face vertex by \texttt{2-extrusion}.
  Once again, the deletion \texttt{DELETE 3-insert} does not modify the orbits.
  %The application \texttt{1-square} has no topological parameter and, thus, is identically reevaluated. Once again, its node $n_5$ only creates dart $5$, and both face vertex $\orbit{1}(5)$ and  volume face $\orbit{0,1}(5)$ are identically reevaluated. Therefore, the orbit level of the reevaluation DAG references dart $5$.
  %\item[Level 2] The added application \texttt{ADD1-insert} does not modify the face vertex $\orbit{1}(5)$ but modifies the volume face $\orbit{0,1}(5)$. 
  %So, the event level 
  %respectively
  %contains \texttt{NOMODIF} and \texttt{MODIFICATION}. Since the application does not match dart $5$, it is kept in the object, thus, the orbit level 
  %still
  %references dart $5$ and contains the same orbits $\orbit{1}$ and $\orbit{0,1}$.
  %\item[Level 3] The \texttt{2-extrusion}'s application is reevaluated identically. The rule's hook $n1$ matches dart $5$ and the tracking allows finding out that the dart created by $n6$, as a copy of dart $5$, is the dart $59$. Therefore, the event level contains two \texttt{CREATION} and the orbit level references dart $59$.
  %\item[Level 4] Since \texttt{3-insert} has been deleted, the application is not reevaluated. Consequently, neither the volume face vertex $\orbit{1}(59)$ nor the volume face $\orbit{0,1}(59)$ are modified.
  %Therefore, 
  %The event level contains two \texttt{NOMODIF} and the orbit level keeps $59$ as a reference and still contains the orbits $\orbit{1}$ and $\orbit{0,1}$. 
  \item[Level 5] The fifth application \texttt{4-extrusion} of the extrusion rule (Fig.~\ref{fig:rule-face-extrusion}) matches the volume face $\orbit{0,1}$ with its hook $n1$ and creates a copy of dart $59$ %being the 
  as dart $63$. As in the initial evaluation, it creates a volume vertex and a face. %Therefore 
  Then, the event level contains two \texttt{CREATION}. The orbit level references dart $63$ and contains both orbits $\orbit{1,2}$ and $@1.@2.\orbit{0,1}$.
  \item[Level 6] Finally, while the tracking of \texttt{5-collapse} shows that the application keeps merging the volume vertices incident to the matched face, it preserves dart $63$ which is matched by node $n2$. Therefore, the event level contains one \texttt{MERGE} and the orbit level references dart $63$ and contains the volume vertex $\orbit{1,2}$.
\end{description}

%Therefore, 
At last, $PN_6$'s concrete name is dart $63$ (cf Fig.~\ref{fig:6-reeval}). %

These two examples here are quite straightforward as there was only one possible candidate dart each time. %
However, in some com\-plex specifications, there can be more than a single dart to choose bet\-ween.

\subsection{Parameter matching strategies}
%\david{L'édition de la specification paramétrique entraine un dag différent au rejeu (avec des event level en plus ou en moins et des des branches en plus ou en moins). 
%PAr exemple la scission d'une orbite qui apparaissait dans le dag initial peut disparaitre au rejeu, une fusion peut être ajoutée, etc. Plusieurs stratégies de matching sont alors possible en fonction du contexte d'application. 
%Cela peut être illustré sur l'exemple de PN7 qui permet de désigner la face qui doit être colorée (fig). 
%L'ajout de ADD-insert au rejeu entraine indirectement la sission de la face qui devait être colorée et l'ajout d'une branche suppléméntaire dans le dag de rejeu...et tu enchaines sur le déroulement de la construction de PN7 : Level 1... Bon c'est à la louche, tu peux reprendre des bouts de ce que tu as déjà écrit.}

The editing of the parametric specification leads to having a dif\-fe\-rent DAG at reevaluation (with event levels and/or branches being added, deleted or both). %
For example, an orbit split present in the evaluation DAG may disappear during the reevaluation, a merging can be added and so on. %
Several matching strategies can then be %envisioned
considered depending on the application's context. %
This can be il\-lus\-tra\-ted with $PN_7$'s example which designates the face that must be coloured (Fig.~\ref{fig:8-eval}). %
The addition of \texttt{ADD1-insert} application at reevaluation splits the origin of designated face, resulting in the addition of a branch in the reevaluation DAG. 

%A topological parameter can be duplicated by an added operation %splitting an orbit. %
%such as an orbit split. %
%A strategy could be that, 
%One could set up a strategy such that,
%during the reevaluation, the application should be applied on as many parameters as there are suborbits resulting from an orbit being split.

\begin{figure}[ht]
    \centering
    \includegraphics[width=\columnwidth]{pn7-reval}
    \caption{$PN_7$ reevaluation DAG}
    \label{fig:pn7-reval-dag}
\end{figure}

%For example, let's consider the topological parameter for \texttt{8-colour}, $PN_7$, being the left face resulting from the triangulation (Fig.~\ref{fig:8-eval}). %
Let's work through $PN_7$'s reevaluation DAG shown in Fig.~\ref{fig:pn7-reval-dag}.%
\begin{description}
  \item[Level 1] As seen previously, \texttt{1-square} creates the tracked orbits traced in the evaluation DAG.
  %  The first application which creates a square face is reevaluated identically and creates a face edge $\orbit{0}$ and a dart orbit $\orbit{}$. %
  %The node $n_3$ creates a single dart $3$. The related event level contains two \texttt{CREATION}; the orbit level references dart $3$ and contains two orbits $\orbit{0}$ and $\orbit{}$.
  \item[Level 2] The added application \texttt{ADD1-insert} matches dart $3$. 
  Since the vertex insertion rule (Fig.~\ref{fig:rule-insert-vertex}) splits the volume face edge $\orbit{0}(3)$, its history is also split and there are two con\-cre\-te names to consider. It follows that the event level contains two \texttt{SPLIT}, one for each volume face edge suborbit, and two \texttt{NOMODIF}, one for each dart $3$ and $4$. 
  %The left orbit level references dart $3$ and contains two orbit $\orbit{0}$ and $\orbit{}$. The right orbit level references dart $4$ and contains the same two orbits $\orbit{0}$ and $\orbit{}$.
  \item[Level 3] The third application \texttt{2-extrusion} extrudes the square face into a cube.
  The extrusion rule (Fig.~\ref{fig:rule-face-extrusion}) matches dart $3$ with its hook $n_1$. It extrudes the face edge $\orbit{0}(3)$ into a volume face, the dart $\orbit{}(3)$ into a face edge and its node $n_2$ creates the dart $37$ as a copy of dart $3$. The same goes for the history on the right. The rule matches dart $4$, it creates the same two orbits and a dart $52$ as a copy of dart $4$. %
  Therefore, both event levels contain two \texttt{CREATION}. Both orbit levels contain orbits $\orbit{0,1}$ and $\orbit{0}$. %
  \item[Level 4] The last application, \texttt{5-triangulation} triangulates the faces designated by the level above. 
  Both darts $37$ and $52$ are matched and preserved by the rule. %, two darts are referenced: dart $37$ on the left orbit level and dart $52$ on the right orbit level. % 
\end{description}

%The reevaluation DAG of $PN_7$ designates two topological parameters from only one persistent name. 
%\david{Cela m'embête de dire qu'il y a 2 paramètres topo. A mon sens il n'y en a qu'un, celui de l'opération (et 1 paramètre topo correspond à 1 PN). Par contre un PN peut être matché sur différentes entités ou orbites au rejeu. Ce terme matché est bien celui que l'on utilise dans la littérature pour parler de ce process. Utilise plutôt cette terminologie} 

%\david{j'ai réorganisé ce que tu disais ainsi : 
Upon reevaluation, $PN_7$'s DAG matches to different darts. An option would be to colour only one face (either $\orbit{0,1}(37)$ or $\orbit{0,1}(52)$). %
Another option would be to apply \texttt{8-colour} two times, one for dart $37$ and one for dart $52$ (as shown in Fig.~\ref{fig:8-reeval} which represent our default strategy). %qui représente la stratégie appliquée par défaut).
%With all the informations (events that occurred on trace and origin orbits) stored within evaluation and reevaluation DAGs,
%Dans la mesure où toutes les informations (évènements survenus sur les orbites traces et origines) sont stockées dans les DAG du jeu et du rejeu
%it is thus possible to show  %
In case the reevaluation DAG has two or more leaves, it shows all the possible entities that can be matched for a specific persistent name and to set a strategy up.
Ei\-ther way, such a strategy allows users to have a choice and best fit their modeling intents, depending on the application's context. %en fonction du contexte d'apllication.
%}

%Following $PN_7$'s reevaluation DAG, two darts can be matched for this persistent name. Therefore, $PN_7$'s possible concrete names are $g=37$ and $h=52$ (cf. Fig.~\ref{fig:7-reeval} and \ref{fig:8-reeval}). 

%Since a reevaluation DAG shows all the possible entities that can be matched for a specific persistent name, it is, thus, possible to set strategies up. Upon reevaluation, $PN_7$'s DAG matches to different darts. An option would be to colour only one face (either $\orbit{0,1}(37)$ or $\orbit{0,1}(52)$). Another option would be to apply \texttt{8-colour} two times, one for dart $37$ and one for dart $52$ (as shown in Fig.~\ref{fig:8-reeval} by default). Either way, such a strategy allows users to have a choice and best fit their modeling intents. %

%\david{Ajouter ici ou ailleurs l'idée suivante qui doit être soulignée :} 
%Note that at level 2, the face edge $\orbit{0}(3)$ is split. %
%This is the origin of the volume face created by \texttt{2-extrusion}, which will be split by \texttt{5-triangulation}, and one of its subfaces will ultimately be colored. %


%\david{Je voulais parler de ça au moment où l'on aborde PN7. Comme on en parle ici, je me suis dit on le met là, mais en relisant je me dis que cette idée serait peut être mieux dans la conclusion, parce qu'elle n'a pas vraiment à voir avec la "stratégie" en elle même.}

%During the initial evaluation (Fig.~\ref{fig:evaluation}), the \texttt{5-triangulation}'s application triangulates only one volume face. %
%However, adding \texttt{ADD1-insert} at reevaluation splits the face edge $\orbit{0}$-orbit being extruded. %
%Following the strategy above, the reevaluation DAG branches after the application of \texttt{ADD1-insert}. %
%From now on, the \texttt{2-extrusion} extrudes each subsequent face edges $\orbit{0}$-orbit into a volume face. %
%Then, each of these faces is triangulated by the \texttt{5-triangulation} application. %
%Finally, the DAG yields the darts $37$ and $52$ as the topological parameters. %
%Applying \texttt{8-colour} on both of these darts leads to having two volume faces ($\orbit{0,1}(37)$ and $\orbit{0,1}(52)$) being coloured in blue (Fig.~\ref{fig:8-reeval}) instead of only one. %

%A rule's application can also lead to the deletion of a topological parameter. %
%When the whole orbit incident to the deleted dart is also deleted, next applications depending on an evolution of said orbit must be cancelled to prevent the reevaluation from failure. %
%When only a part of the orbit is deleted, a viable strategy would be to pick another dart from this orbit as a substitute. % \david{utile uniquement si le brin désigné fait partie du morceau supprimé.}
%Since the application of a rule on two different darts belonging to the same orbit yield the same results, this strategy would make it possible to %pursue 
%continue the reevaluation until %it is done. %
%its completion. % \david{Pas sûr que cela vaille la peine de parler de ce cas où il faut trouver un autre représentant.}
%\david{perso, je virerai tout ce dernier paragraphe... en plus ca fait gagner un peu de place ;-)}